% page 199 du fac-simile (image p-198 du scan)
% bloc 175.3 x 259.3 mm ; marge gauche 26.7 mm
\pgc{5.080mm}{0.000mm}{
\l{}
\l{}
\l{}
\l{\cel{59}191}
\l{}
\l{}
\l{\sou{generatoro}. (\sou{tekn.)}\cel{2}Termino teknikala, kelke sinonima ye \textquotedbl{}kaldiero\textquotedbl{},}
\l{\cel{3}qua tante plu uzesas, quante plu la tipi difereskas de la tanki}
\l{\cel{3}cilindra, tre kapaca, quin onu nomizas \textquotedbl{}kaldieri\textquotedbl{} ; la genera-}
\l{\cel{3}tori esas generale aquo-tuboza e rezistas a presi normala ye}
\l{\cel{3}80 kilogrami sur omna centimetro quadrata. - DEF.}
\l{}
\l{\sou{genezo}. Produktado di la enti organizita. - DEFIS.}
\l{}
\l{\sou{genio}. I. Ento supernatura, benigna o maligna, qua (krede da la}
\l{\cel{3}pagani) direktis la fato di la individui, familii, nacioni. -}
\l{\cel{3}II. Ento supernatura qua (en la fe-rakonti) dotesis ye povo}
\l{\cel{3}magiala. - III. Ento alegoriala qua enkorpigas (en poezio,}
\l{\cel{3}skultarto, piktarto) ideo abstraktita. - IV. Talento naturala.}
\l{\cel{3}- V. Kapableso supera quan Naturo grantis a la spiriti kreera.}
\l{\cel{3}- DEFIRS.}
\l{}
\l{\sou{genitar}. I(\sou{trans.}) Produktar per la fekundigo di la femino. -}
\l{\cel{3}II. (\sou{geom.}) Produktar per movar su (ex.: \textquotedbl{}Sfero genitesas per}
\l{\cel{3}la rotaco di mi-cirklo cirkum sua diametro\textquotedbl{}). - DEFIRS.}
\l{}
\l{\sou{genitivo}. (\sou{gramat.)}\cel{2}Kazo di vorto deklinebla, qua expresas lua}
\l{\cel{3}relato ad altra vorto qua indikas kozo qua apartenas ad ita.}
\l{\cel{3}- DEFIS.}
\l{}
\l{\sou{genitori}. Gepatri.}
\l{}
\l{\sou{genro}. (\sou{gram.}) Formo specala, uzata da ula lingui, por distingar}
\l{\cel{3}la enti maskula, la enti femina e la kozi sensexua.}
\l{}
\l{\sou{gento}. I. (\sou{che la Romani antiqua)}. Familio o grupo de familii qua}
\l{\cel{3}havis origino komuna, patrica. - II. La sama-rasani. - L.}
\l{}
\l{\sou{genuo}. (\sou{anat.})\cel{2}Artiko di la gambo e di la kruro, ye la parto}
\l{\cel{3}avana. - EFI.}
\l{}
\l{\sou{geocentrala}. (\sou{astron.})\cel{2}Mezurata relate Tero quan onu konsideras}
\l{\cel{3}kom centro di la cie lo-sfero. - DEF.}
\l{}
\l{\sou{geodezio}. Cienco pri la mezuro di la Terglobo o di lua parti.DEFIRS.}
\l{}
\l{\sou{geognozio}. Cienco pri la kompozeso mineralogiala di la Terglobo.}
\l{}
\l{\sou{geografo}. La persono qua su okupas pri geografio. - DEFIRS.}
\l{}
\l{geografio. Cienco di qua la studiajo esas la deskripto di la}
\l{\cel{3}Teresurfaco. - DEFIRS.}
\l{}
\l{geologo. La persono qua su okupas pri geologio. - DEFIRS.}
\l{}
\l{geologio. Cienco qua havas kom studiajo la strati ye qui konsistas}
\l{\cel{3}la Ter-kortico,lia estado nuna e lia formaceso. - DEFIRS.}
\l{}
\l{\cel{1}\sou{geomanciar}. (\sou{netrans}.)Divinar lo futura segun la figuri quin forma-}
\l{\cel{3}cas manuedo de tero jetita sen-vice. - DEF.}
}
